\subsection{Parameter convergence: retaining the adaptive-direction defect}
\label{app:convergence}

Unlike the routing variable, the parameter is updated in Euclidean space.
The appropriate smoothness inequality is always valid for an Adam direction;
what is not automatic is a positive conditional alignment between that
direction and the population gradient. We make this distinction explicit.
Let $f$ be the objective under analysis, either empirical risk on a fixed
training set or population risk. Those objectives must not be interchanged
when claiming unbiasedness.

\begin{assumption}[Smooth objective and a predictable comparison metric]
\label{ass:parameter}
$f$ is $L$-smooth and bounded below by $f_\star$.
$g_t=\nabla f(W_t)$ and
$\mathbb E_tG_t=g_t$,
$\mathbb E_t\|G_t-g_t\|^2\le\sigma_G^2$.
An $\mathcal F_t$-measurable self-adjoint linear map
$\mathcal B_t$ on matrix parameters satisfies
\begin{equation}
 a\|X\|^2\le\langle X,\mathcal B_tX\rangle,
 \qquad\|\mathcal B_tX\|\le b\|X\|,
 \qquad 0<a\le b<\infty.\label{eq:metric-bounds}
\end{equation}
The implemented parameter direction is decomposed \emph{exactly} as
\begin{equation}
 D_t=\mathcal B_tG_t+e_t,\qquad
 \mathbb E_t\|e_t\|^2\le E_t^2.\label{eq:parameter-defect}
\end{equation}
The learning rates are deterministic, positive, and
$\eta_t\le a/(4Lb^2)$.
\end{assumption}

This assumption isolates a measurable error budget rather than presuming
that Adam is a fresh gradient step. For example, row $i$ can use the
predictable operator
\begin{equation}
 \mathcal B_{t,i}=U_t A_{t,i}U_t^\top+\gamma_t(I-P_t),
 \quad aI_r\preceq A_{t,i}\preceq bI_r,
 \quad a\le\gamma_t\le b.\label{eq:routed-metric}
\end{equation}
Here $A_{t,i}$ may be a clipped inverse square root of the \emph{previous}
second moment. It is not the denominator formed from the current batch.
The actual direction minus this comparison direction defines $e_t$; this
difference includes moment lag, same-batch denominator dependence,
diagonal-transport approximation, and sampled residual normalization. With
weight decay, its direction is included in $e_t$ when $f$ is unregularized;
alternatively one can formulate the objective and comparison update for the
regularized problem. A decoupled decay term must not simply disappear from
the proof.

\begin{theorem}[Nonconvex convergence with explicit adaptive defects]
\label{thm:convergence}
Under Assumption~\ref{ass:parameter}, let
$S_T=\sum_{t<T}\eta_t$ and draw $\tau\in\{0,\ldots,T-1\}$ with
$\Pr(\tau=t)=\eta_t/S_T$. For
$W_{t+1}=W_t-\eta_tD_t$,
\begin{align}
 \mathbb E\|\nabla f(W_\tau)\|^2
 &\le\frac{4(f(W_0)-f_\star)}{aS_T}
       +\frac{2}{a^2S_T}\sum_{t<T}\eta_t\mathbb E E_t^2\notag\\
 &\quad+\frac{4L}{aS_T}\sum_{t<T}\eta_t^2
                         (b^2\sigma_G^2+\mathbb EE_t^2).
 \label{eq:parameter-convergence}
\end{align}
For $\eta_t=\eta_0/\sqrt T$, the noise and initial-suboptimality terms
are $O(T^{-1/2})$. The same stationarity rate follows if
$T^{-1}\sum_{t<T}\mathbb EE_t^2=O(T^{-1/2})$.
For an arbitrary actual Adam run the displayed defect terms remain in the
bound; unbiased raw gradients alone do not make them vanish.
\end{theorem}

\begin{proof}
Predictability is needed at precisely the following step:
\begin{equation}
 \mathbb E_t\langle g_t,\mathcal B_tG_t\rangle
 =\langle g_t,\mathcal B_tg_t\rangle\ge a\|g_t\|^2.
\end{equation}
Using Cauchy--Schwarz and Young on the remaining term gives
\begin{align}
 \mathbb E_t\langle g_t,D_t\rangle
 &\ge a\|g_t\|^2-\|g_t\|\sqrt{\mathbb E_t\|e_t\|^2}\\
 &\ge\tfrac a2\|g_t\|^2-\tfrac{E_t^2}{2a}.
 \label{eq:alignment-bound}
\end{align}
Furthermore,
\begin{equation}
 \mathbb E_t\|D_t\|^2
 \le2b^2(\|g_t\|^2+\sigma_G^2)+2E_t^2.
 \label{eq:direction-bound}
\end{equation}
Smoothness applied to the actual update yields
\begin{align}
 \mathbb E_tf(W_{t+1})
 &\le f(W_t)-\eta_t\mathbb E_t\langle g_t,D_t\rangle
                   +\tfrac L2\eta_t^2\mathbb E_t\|D_t\|^2\\
 &\le f(W_t)-\tfrac a4\eta_t\|g_t\|^2
      +\tfrac{\eta_tE_t^2}{2a}
      +L\eta_t^2(b^2\sigma_G^2+E_t^2).
\end{align}
The last line uses $L\eta_tb^2\le a/4$.
Taking expectation, summing, and using $f(W_T)\ge f_\star$ proves
\eqref{eq:parameter-convergence}. Weighted output sampling gives its
left side exactly. Substituting the constant-in-horizon step proves the
rate statements.
\end{proof}

\begin{corollary}[A deterministic guarantee for the implemented Adam branch]
\label{cor:deterministic-adam}
Consider full-gradient GrassRoute without weight decay, with zero initial
first moments, Procrustes first-moment alignment, and fixed
$\beta_1\in[0,1)$. Assume $\|\nabla f(W_t)\|\le G_0$, the
bias-corrected routed inverse denominators have eigenvalues in $[a,b]$,
and the actual residual gain is in the same interval, $a>0$.
Suppose $\|D_t\|\le D_0$, $\|H_t\|\le H_0$, and the horizon-dependent
steps $\eta,\rho$ are constant within the run and satisfy the small-step
conditions above. Then the actual first-moment contribution to the
direction defect satisfies
\begin{equation}
 E_t\le \frac{b\beta_1}{1-\beta_1}
                 (LD_0\eta+2G_0H_0\rho).\label{eq:deterministic-lag}
\end{equation}
Consequently $\eta,\rho=\Theta(T^{-1/2})$ give
\begin{equation}
 \frac1T\sum_{t<T}\|\nabla f(W_t)\|^2=O(T^{-1/2}).
\end{equation}
The second moments may depend on the current full gradient: no
predictability assumption is needed in this deterministic specialization.
This is a guarantee for the actual Adam first-moment update under bounded
gains, not a claim for arbitrary stochastic Adam.
\end{corollary}
\begin{proof}
Let $m_t=M_tU_t^\top$ denote the ambient first moment immediately after
its EMA update and before transport. Norm preservation of Procrustes
alignment and $\|G_tU_t\|\le G_0$ give
$\|M_t\|\le(1-\beta_1^t)G_0$ by induction.
Write $\delta_{t-1}=d_{\mathrm{Gr}}(U_t,U_{t-1})$.
The aligned frame mismatch has operator norm
$2\sin(\theta_{\max}/2)\le\delta_{t-1}$, so the preceding moment's
ambient transport error is at most
$(1-\beta_1^{t-1})G_0\delta_{t-1}$.
Also
\begin{equation}
 \|g_tP_t-g_{t-1}P_{t-1}\|
 \le LD_0\eta+G_0\delta_{t-1}.
\end{equation}
For $r_t=\|m_t-(1-\beta_1^t)g_tP_t\|$, subtract the target from the
EMA recursion and use these two inequalities:
\begin{equation}
 r_t\le\beta_1r_{t-1}
 +\beta_1(1-\beta_1^{t-1})(LD_0\eta+2G_0\delta_{t-1}).
\end{equation}
Lemma~\ref{lem:qr} gives $\delta_{t-1}\le\rho H_0$.
Starting from $r_1=0$, summing the geometric recursion and dividing by
$1-\beta_1^t$ bounds the bias-corrected ambient moment error by
$\beta_1(LD_0\eta+2G_0H_0\rho)/(1-\beta_1)$.
Lifting is isometric in the routed coordinates; multiplication by the
current inverse denominator has norm at most $b$. This proves
\eqref{eq:deterministic-lag}. Choose the comparison metric using that
current denominator and the actual residual gain. Pathwise positive
definiteness gives the alignment step in the proof of
Theorem~\ref{thm:convergence}, without a stochastic conditional
expectation. The resulting bound with $\sigma_G=0$ has terms
$O((\eta T)^{-1})+O((\eta+\rho)^2)$, proving the stated rate.
\end{proof}

The gain restrictions have concrete interpretations. A positive Adam
denominator offset bounds inverse gains above; bounded corrected second
moments bound them below. With normalized covariance samples, tangency
projection is nonexpansive and $\|\widehat Q_t\|\le1$, so zero-initialized
Grassmann EMA has $H_0\le1$. A positive residual clipping lower bound
ensures residual coercivity. If the implementation instead permits its
gain to approach zero, that uniform lower-bound assumption has not been
verified. The result therefore applies to its bounded-gain regime or a
specified positive-floor configuration, not automatically to every run.

\paragraph{A rigorous specialization and a counterexample.}
The theorem applies without an adaptive defect to a predictable, bounded
routed preconditioner with a positive residual branch. In deterministic
training it also applies directly to a current-gradient positive definite
preconditioner, using pathwise alignment instead of predictability and the
corresponding elementary descent argument. These are genuine sufficient
conditions. Standard Adam with exponential moments and a same-batch
denominator is more delicate \cite{reddi2018adam,chen2019adamtype}.
For a simple illustration, take a scalar unbiased gradient with
$G=10$ with probability $0.1$ and $G=-1$ otherwise. Then
$\mathbb EG=0.1>0$, but a fresh zero-history Adam step with negligible
$\varepsilon$ has $\mathbb E\operatorname{sign}(G)=-0.8$.
The expected direction is uphill for this positive population gradient.
Unbiasedness of $G$ therefore cannot justify the cancellation used above
for a denominator depending on $G$. Momentum can also temporarily point
uphill after a gradient sign change. These examples explain why
\eqref{eq:parameter-defect} is retained for the implemented algorithm.

\subsection{How tracking, moment transport, and rank tails enter convergence}

The theorem can be connected to the routing analysis without conflating
their covariances. Let $\mathcal B_t^\star$ be an ideal predictable metric
constructed in a chosen oracle subspace, with the same positive bounds as
\eqref{eq:metric-bounds}, and write
\begin{equation}
 D_t=\mathcal B_t^\star G_t+
       (\mathcal B_t-\mathcal B_t^\star)G_t+e_t^{\rm state}.
\end{equation}
If, after consistent coordinate alignment,
\begin{equation}
 \|\mathcal B_t-\mathcal B_t^\star\|_{\rm op}
 \le K_P\|P_t-P_t^\star\|_{\rm op}+d_t^{\rm metric},
 \label{eq:metric-perturbation}
\end{equation}
and $\mathbb E_t\|G_t\|^2\le G_0^2$, then
\begin{align}
 \mathbb E_t\|e_t\|^2
 \le 3G_0^2K_P^2\sin^2\theta_{\max,t}
       +3G_0^2(d_t^{\rm metric})^2
       +3\mathbb E_t\|e_t^{\rm state}\|^2.
 \label{eq:tracking-to-convergence}
\end{align}
The projector perturbation term is geometric; the metric term is separate
because different coordinatewise diagonal preconditioners can act
differently in the \emph{same} subspace. Equation
\eqref{eq:metric-perturbation} is an extra operator-consistency condition,
not a consequence of matching projectors.
If its oracle is the leading eigenspace of the normalized covariance,
Theorem~\ref{thm:tracking} controls the first term. If the oracle instead
uses $S_t$, normalization mismatch adds another projector perturbation,
which can be bounded by Davis--Kahan only when the corresponding raw
covariance gap is available.

\begin{proposition}[Residual preservation removes a mandatory rank-tail floor]
\label{prop:no-tail-floor}
For the metric \eqref{eq:routed-metric} with $\gamma_t\ge a>0$,
\eqref{eq:parameter-convergence} contains no obligatory
$\overline{\tau_r}$ term. Even rank zero reduces to a scaled full-gradient
method and can converge. If the residual is deliberately removed,
the missing population-gradient energy can be bounded by
\begin{equation}
 \|g_t(I-P_t)\|^2
 \le\mathbb E_t\|G_t(I-P_t)\|^2
 =\tau_r(S_t)+2h_{S_t}(U_t).\label{eq:dropped-tail}
\end{equation}
Consequently a rank-tail floor is a valid upper-bound contribution for a
\emph{projection-only} analysis, not an impossibility theorem for full
GrassRoute.
\end{proposition}
\begin{proof}
Orthogonality of the two branches gives
\begin{equation}
 \langle X,\mathcal B_tX\rangle
 =\sum_i\langle X_iU_t,A_{t,i}X_iU_t\rangle
       +\gamma_t\|X(I-P_t)\|^2\ge a\|X\|^2.
\end{equation}
The comparison direction is therefore coercive on the full parameter
space, regardless of $r$ and its covariance tail. Jensen's inequality for
the projected gradient, followed by Lemma~\ref{lem:decomposition}, proves
\eqref{eq:dropped-tail} when the residual is omitted.
\end{proof}

For a scalar routed gain $a_t$ and residual gain $\gamma_t$, there is an
even simpler identity:
\begin{equation}
 D_t=a_tG_t+(\gamma_t-a_t)G_t(I-P_t),\qquad
 \mathbb E_t\|D_t-a_tG_t\|^2
 =(\gamma_t-a_t)^2\{\tau_r(S_t)+2h_{S_t}(U_t)\}.
 \label{eq:tail-gain}
\end{equation}
Here a heavy tail increases the discrepancy from an isotropic reference
only when the branch gains differ. Large residual energy can hurt an
optimization mechanism empirically, but cannot by itself establish an
irreducible stationarity error for a method that retains it.

\paragraph{Bounding moment-induced defects.}
The scalar map $a\mapsto1/(a+\varepsilon)$ on $a\ge0$ is
$\varepsilon^{-2}$-Lipschitz. Working with the root-amplitude state
$A=\sqrt V$ therefore avoids an unjustified derivative bound for
$V\mapsto V^{-1/2}$ near zero. For bounded $\|M\|\le M_0$,
\begin{equation}
 \left\|\frac{M}{A+\varepsilon}-
              \frac{M'}{A'+\varepsilon}\right\|
 \le\varepsilon^{-1}\|M-M'\|
        +M_0\varepsilon^{-2}\|A-A'\|.\label{eq:adam-lipschitz}
\end{equation}
Bias correction contributes its explicit finite factors
$(1-\beta_1^t)^{-1}$ and $(1-\beta_2^t)^{-1/2}$.
For a fixed-coordinate momentum EMA,
\begin{equation}
 M_t=(1-\beta_1)\sum_{k=0}^{t-1}\beta_1^k Z_{t-k},
\end{equation}
and thus
\begin{equation}
 \|M_t-(1-\beta_1^t)Z_t\|
 \le(1-\beta_1)\sum_{k<t}\beta_1^k\|Z_{t-k}-Z_t\|.
\end{equation}
If consecutive aligned coordinates change by at most $D_Z$, the
right side is at most $\beta_1D_Z/(1-\beta_1)$.
Moving coordinates additionally contribute the first-moment projection
error and second-moment covariance error in Lemma~\ref{lem:moments}.
Combining these bounds with \eqref{eq:adam-lipschitz} gives explicit
state contributions to $E_t$ once their magnitudes are controlled.
It does not establish that those magnitudes vanish in arbitrary neural
training. Resetting a denominator may increase its reciprocal toward
$1/\varepsilon$ and amplify precisely these terms.

\begin{proposition}[Gradient norm matching does not balance branch steps]
\label{prop:branch-scale}
Consider prescribed nonzero gradient histories $G_1,\ldots,G_t$ and
$cG_1,\ldots,cG_t$, $c>0$, with identical initial basis and zero moments.
In the zero-offset idealization of normalized covariance routing, diagonal
moment transport, and residual norm matching, the two basis histories agree,
their first and second moments scale by $c$ and $c^2$, respectively, and
\begin{equation}
 D_{t,c}^{\parallel}=D_t^{\parallel},\qquad
 D_{t,c}^{\perp}=cD_t^{\perp}.
 \label{eq:branch-scaling}
\end{equation}
The statement applies to coordinates with nonzero second moments, and a
zero residual remains zero. For a fixed branch learning-rate ratio, the
routed-to-residual update norm ratio consequently scales as $1/c$.
With positive offsets this is an approximate statement in the regime where
the corresponding scaled norms dominate the offsets.
\end{proposition}
\begin{proof}
Normalized covariance samples are invariant under positive scaling of $G$,
so the tangent directions, subspace momentum, QR bases, and Procrustes
rotations are the same. Induction in the moment updates gives
$M_{t,c}=cM_t$ and $V_{t,c}=c^2V_t$; diagonal covariance transport is
homogeneous and preserves these relations. Bias correction has identical
factors. Thus $c\widehat M_t/\sqrt{c^2\widehat V_t}
=\widehat M_t/\sqrt{\widehat V_t}$ coordinatewise.
The norm ratio $\|G\|/\|G^\perp\|$, including its fixed clipping limits,
is scale invariant, whereas its product with $G^\perp$ scales by $c$.
This proves \eqref{eq:branch-scaling}.
\end{proof}

For an unclipped, negligible-offset residual rule,
$\|D_t^\perp\|\simeq\|G_t\|$, while the routed Adam direction can
remain approximately scale invariant. A small gradient therefore does not
make the two branch \emph{update} norms comparable. This provides an
additional finite-budget failure mechanism beyond subspace tracking:
residual directions can be updated too slowly relative to the adaptive
branch even while their gradient is preserved. In particular, identical
normalized eigengaps, captured-energy fractions, and principal angles do
not determine the entire parameter dynamics. This observation does not
establish a nonzero asymptotic error floor, nor identify the sole cause of
any benchmark outcome; measuring actual branch step norms would be needed
for that causal claim. It is consistent with the importance of recovery
scaling studied in related optimizers \cite{chen2024fira,rajabi2025subtrackpp}.

\paragraph{Iteration and time complexity.}
When the theorem's defect condition holds, a target
$\mathbb E\|\nabla f\|^2\le\epsilon$ requires
$O(\epsilon^{-2})$ iterations for the stated stochastic schedule, with
constants depending on $a,b,L,\sigma_G$. This is a stationarity statement,
not global optimality or a test-accuracy guarantee. Routing costs
$O(mdr)$ for $GU$, $G^\top(GU)$, and lifting, plus $O(dr^2+r^3)$ for QR
and the small Procrustes polar factor. No $d\times d$ covariance need be
formed by the optimizer. Persistent routed tensors $M,V,U,H$ use exactly
$2mr+2dr$ scalar entries, before fallback states. Temporary full gradients,
activations, and parameter memory are additional. A step-complexity
advantage must therefore be multiplied by actual per-step time before
making a wall-clock claim.
